% Direct-authority diplomatic French for printed p.134.
\SourcePage{134}

\LayerTarget{tohoku:prop:1.8:proof:p2}
\noindent $=w_{ij}(A_i)$, soit $N$ le sous-truc de $\mathrm{S}$ borne supérieure des
$N_{ij}$ (qui existe en vertu de l'axiome AB 3)), soit
$\overline{\mathbf A}=\mathrm{S}/N$, soit $u_i:A_i\to\overline{\mathbf A}$ le morphisme
induit par le morphisme canonique $\mathrm{S}\to \mathrm{S}/N$, on constate aussitôt
que $(\mathbf A,u_i)$ est une limite inductive de $\mathbf A$. Nous laissons
au lecteur la démonstration des autres assertions de la proposition 1.8.,
démonstration évidemment bien connue.

\LayerTarget{tohoku:section:1.9}
\textbf{1.9. Générateurs et cogénérateurs.} Soit $\mathbf{C}$ une catégorie,
et soit $(U_i)_{i\in I}$ une famille d'objets de $\mathbf{C}$. On dit que c'est une
\textit{famille de générateurs} de $\mathbf{C}$ si pour tout objet $A\in \mathbf{C}$ et tout
sous-truc $B\neq A$, on peut trouver un $i\in I$ et un morphisme
$u:U_i\to A$ qui ne provienne pas d'un morphisme de $U_i$ dans
$B$. Alors pour tout $A\in \mathbf{C}$, les sous-trucs de $A$ forment un
\textit{ensemble} : en effet, un sous-truc $B$ de $A$ est complètement
déterminé par l'ensemble des morphismes d'objets $U_i$ dans $A$ qui proviennent
d'un morphisme de $U_i$ dans $B$. On dit qu'un objet $U\in \mathbf{C}$ est un
\textit{générateur} de $\mathbf{C}$ si la famille $\{U\}$ est une famille de
générateurs.

\LayerTarget{tohoku:prop:1.9.1}
\textsc{Proposition 1.9.1.} \textit{Supposons que $\mathbf{C}$ soit une
catégorie abélienne satisfaisant à l'axiome \textup{AB 3)} (existence de sommes directes
infinie), soit $(U_i)_{i\in I}$ une famille d'objets de $\mathbf{C}$,
$U=\oplus U_i$ sa somme directe. Les conditions suivantes sont
équivalentes :}

\textit{a) $(U_i)$ est une famille de générateurs de $\mathbf{C}$.}

\textit{b) $U$ est un générateur de $\mathbf{C}$.}

\textit{c) Tout $A\in \mathbf{C}$ est isomorphe à un quotient d'une somme
directe $U^{(I)}$ d'objets tous identiques à $U$.}

\LayerTarget{tohoku:prop:1.9.1:proof}
L'équivalence de a) et b) est une conséquence à peu près immédiate de la
définition. b) implique c), car il suffit de prendre pour $I$ l'ensemble
$\mathrm{Hom}(U,A)$ et de considérer le morphisme de $U^{(I)}$ dans $A$ dont
la composante suivant chaque $u\in I$ est $u$ lui-même : l'image $B$ de ce
morphisme est $A$, puisque autrement il existerait un
$u\in\mathrm{Hom}(U,A)=I$ tel que $u(A)\not\subset B$, ce qui serait absurde ;
donc $A$ est isomorphe à un quotient de $U^{(I)}$. c) implique b), car il est
immédiat que si $A$ est un quotient de $U^{(I)}$, alors pour tout sous-truc
$B$ de $A$ distinct de $A$, il existe un $i$ tel que l'image canonique dans
$A$ du $i$.ème facteur de $U^{(I)}$ soit $\not\subset B$, d'où un morphisme de
$U$ dans $A$ qui ne provient pas d'un morphisme de $U$ dans $B$. (On notera
que la structure additive de $\mathbf{C}$ n'a pas réellement servi ici).

\LayerTarget{tohoku:section:1.9:examples:p1}
\textsc{Exemples.} Si $\mathbf{C}$ est la catégorie abélienne des modules à
gauche unitaires sur un anneau $U$ avec unité, alors $U$ (considéré comme
module à gauche sur lui-même) est un générateur. Si $\mathbf{C}$ est la catégorie des
faisceaux de groupes abéliens sur un espace topologique fixé $X$, et si pour
tout ouvert $U\subset X$ on désigne par $\mathbf{Z}_U$ le faisceau sur $X$ qui est nul
au dessus de $\complement U$ et identique au faisceau constant des entiers au dessus de
$U$, la famille des $\mathbf{Z}_U$ forme un système de générateurs de $\mathbf{C}$ ; cet exemple
se généralise aussitôt au cas où on se donne sur $X$ un faisceau $\mathbf{O}$ d'anneaux,
et où l'on considère la catégorie des faisceaux de $\mathbf{O}$-modules sur $X$.
D'autres exemples sont contenus dans la proposition suivante :

\TohokuContinuation{135}
